%%% ВЫБОР УСТРОЙСТВА %%%

\IfStrEqCase{\devicecode}{%
{ЮТКБ.656122.611-05.01 РЭ}{%
% код для этого типа
\def\fileExt{dzo1}
\def\AbzazOne{
\item
Алгоритм работы УРОВ выключателя В1 приведен на рисунке \ref{pict_urov:UROV_Q1}. Алгоритмы работы УРОВ выключателей В2, В3, В4 выполнены аналогично. Параметры для настройки УРОВ В1..УРОВ В4 приведены в таблицах \ref{app_set:urov1} - \ref{app_set:urov4}.
}
}%
{ЮТКБ.656122.611-05.02 РЭ}{%
% код для этого типа
\def\fileExt{dzo2}
\def\AbzazOne{
\item
Алгоритм работы УРОВ выключателя В1 приведен на рисунке \ref{pict_urov:UROV_Q1}. Алгоритм работы УРОВ выключател В2 выполнен аналогично. Параметры для настройки УРОВ В1, УРОВ В2 приведены в таблицах \ref{app_set:urov1}, \ref{app_set:urov2}.
}
}%
}
[
\def\fileExt{??}
\def\AbzazOne{??}
]


\color{unidarkgreen}{\subsection{ Устройство резервирования отказа выключателя}}
\color{black}

\begin{enumerate}

%%%%%%%%%%%%%%%%%%%%%%%%%%%%%%%%%%%%%%%%%%%%%%%%%% Общие данные %%%%%%%%%%%%%%%%%%%%%%%%%%%%%%%%%%%%%%%%%%%%%%%%%%
\item \textbf{Общие данные}

\begin{enumerate}

\item
Функция УРОВ предназначена для быстрого и селективного отделения от энергосистемы поврежденного присоединения с отказавшим выключателем в аварийных ситуациях путем отключения смежных выключателей, и с возможностью предварительного повторного действия на отключение отказавшего выключателя.

\item
Логика работы устройства обеспечивает возможность реализации индивидуального УРОВ присоединения как в составе терминала ДЗО (внутренний УРОВ), так и в составе терминалов защит и АУВ присоединений (внешний или распределенный УРОВ). Выбор режима работы УРОВ осуществляется с помощью программного переключателя \textbf{«Режим УРОВ Вn»}. При установленном положении \textbf{«Внеш»} используется схема внешнего УРОВ, реализованная в терминалах защит и АУВ присоединений. При положении \textbf{«Внут»} используется схема внутреннего УРОВ, реализованная в терминале ДЗО.

\end {enumerate}


%%%%%%%%%%%%%%%%%%%%%%%%%%%%%%%%%%%%%%%%%%%%%%%%%% Внутренний УРОВ %%%%%%%%%%%%%%%%%%%%%%%%%%%%%%%%%%%%%%%%%%%%%%%%%%
\item \textbf{Внутренний УРОВ}

\begin{enumerate}

\item
Внутренний УРОВ обеспечивает действие как на собственный выключатель («на себя»), так и на смежные выключатели. Действие на смежные выключатели включает в себя:
\begin{itemize}
\item отключение смежных выключателей через защиты смежных присоединений;
\item останов ВЧ-передатчика защиты;
\item пуск аварийной команды ТО на противоположную сторону ЛЭП;
\item запрет АПВ выключателей.
\end{itemize}

\item
Пуск УРОВ осуществляется при срабатывании внутренних защит на отключение (\textit{«ЛО Вn/ ЛО: Отключение»}), при наличии сигнала пуска УРОВ от внешних защит (\textit{«Пуск УРОВ Вn»}) или сигнала пуска УРОВ НН (\textit{«Пуск УРОВ Bn НН»}).

\item
Для подтверждения пуска УРОВ при КЗ предусмотрен контроль наличия тока через выключатель. Контроль осуществляется с помощью ИО, ток срабатывания которого задается уставкой \textbf{«Iср УРОВ Вn»}.

\item
Для обеспечения надежного пуска УРОВ при кратковременном пропадании сигналов пуска до полного обесточивания отказавшего выключателя предусмотрен подхват от ИО УРОВ, ввод которого осуществляется 
с помощью программного ключа \textbf{«Фикс. УРОВ Вn»}. Сброс пуска произойдет только при возврате ИО УРОВ.

\item
Сигнал \textit{«Пуск УРОВ Вn НН»} предусмотрен для пуска УРОВ от защит, срабатывание которых не сопровождается увеличением тока или сопровождается протеканием малых токов КЗ, недостаточных для срабатывания токовых ИО УРОВ. Например, при пуске УРОВ от УРОВ НН при КЗ на стороне НН (авто)трансформатора за токоограничивающим реактором или при пуске УРОВ от газовой защиты при малой нагрузке или ее отсутствии.
Данный пуск осуществляется с контролем положения выключателя. Фиксация включенного положения выключателя осуществляется по инверсному сигналу \textit{«УРОВ Вn/ УРОВ: В отключен»}.

\item
В устройстве предусмотрена возможность реализации двух принципов выполнения внутреннего УРОВ:
\begin{itemize}
\item с автоматической проверкой исправности выключателя;
\item с дублированным пуском.
\end{itemize}

\item
Для реализации принципа УРОВ с дублированным пуском необходимо задать следующие значения уставок: \textbf{«Режим УРОВ Вn с дубл пуском»} = «Введено», \textbf{«УРОВ Вn на себя»} = «Выведено».
В этом случае УРОВ действует только на отключение смежных выключателей с дополнительным контролем положения «своего» выключателя. Контроль осуществляется по сигналу \textit{«УРОВ Вn/ УРОВ: Готовность ЦО»}.

\item
Для реализации принципа УРОВ с автоматической проверкой исправности выключателя необходимо задать следующие значения уставок: \textbf{«Режим УРОВ Вn с дубл пуском»} = «Выведено», \textbf{«УРОВ Вn на себя»} = «Введено». В этом режиме УРОВ действует на отключение «своего» выключателя при наличии пусковых условий в течение выдержки времени \textbf{«Тср УРОВ Вn на себя»}. Повторное отключение позволяет исключить ложное срабатывание УРОВ на смежные выключатели за счет возврата ИО УРОВ после отключения выключателя. Действие УРОВ «на себя» может быть выполнено с контролем по тока при введенном программном ключе \textbf{«Контр по I УРОВ Вn на себя»}.

\end {enumerate}


%%%%%%%%%%%%%%%%%%%%%%%%%%%%%%%%%%%%%%%%%%%%%%%%%% Внешний УРОВ %%%%%%%%%%%%%%%%%%%%%%%%%%%%%%%%%%%%%%%%%%%%%%%%%%
\item \textbf{Внешний УРОВ}

\begin{enumerate}

\item
Логика внешнего УРОВ предусматривает возможность приема сигналов срабатывания УРОВ от терминалов защит и АУВ присоединений. Для этого используется внешний сигнал \textit{«Ср. УРОВ Вn»}.

\item
С целью повышения надёжности работы и исключения ложных срабатываний реализован токовый контроль, осуществляемый с помощью ИО внутренних УРОВ.

\item
Ввод и вывод функции УРОВ осуществляется либо программным переключателем \textbf{«Режим УРОВ Вn»}, либо посредством сигналов \textit{«УРОВ Вn/ УРОВ: Ремонт»} или \textit{«УРОВ Вn: Вывод»}.

\AbzazOne

\medskip
\begin{figure}[H]
\centering
\includegraphics[width=1\textwidth,height=1\textheight,keepaspectratio]{\fbpath/05.01 ДЗО/661.3331 УРОВ/img/UROV_Q1_\fileExt.pdf}
\captionof{figure}{Функциональная схема логики УРОВ В1}\label{pict_urov:UROV_Q1}
\end{figure}
\medskip

\item
Сигнал срабатывания УРОВ СШ формируется при срабатывании одного из внутренних УРОВ присоединений с учетом текущей фиксации присоединения с отказавшим выключателем или при наличии сигналов срабатывания внешнего централизованного УРОВ \textit{«УРОВ СШ/ УРОВ: Центр. УРОВ»}.

\item
При длительном наличии сигналов пуска и срабатывания УРОВ более выдержки времени \textbf{«Тнеиспр УРОВ»} формируется сигнал неисправности цепи УРОВ.

\item
Алгоритм работы УРОВ СШ приведен на рисунке \ref{pict_urov:UROV_SSH}. Параметры для настройки УРОВ СШ приведены в таблице \ref{app_set:urovsh}.

\medskip
\begin{figure}[H]
\centering
\includegraphics[width=1\textwidth,height=1\textheight,keepaspectratio]{\fbpath/05.01 ДЗО/661.3331 УРОВ/img/UROV_SSH_\fileExt.pdf}
\captionof{figure}{Функциональная схема логики УРОВ СШ}\label{pict_urov:UROV_SSH}
\end{figure}
\medskip

\end {enumerate}

\end{enumerate}